\documentclass[border=12pt]{standalone}
\usepackage[utf8]{inputenc}
\usepackage[T1]{fontenc}
\usepackage[american, siunitx]{circuitikz}

\sisetup{output-decimal-marker={,}}

\ctikzset{
    resistors/scale=0.8,
    diodes/scale=0.8
}

\begin{document}
\begin{circuitikz}[font=\sffamily\small]

    % =========================================================================
    % 1. BLOCO ESQUERDO: ARDUINO UNO R3 (U1)
    % =========================================================================
    \draw[thick, fill=gray!5] (0, -7.0) rectangle (4.0, 7.0);
    \node[align=center, font=\sffamily\bfseries] at (2.0, 0.4) {U1\\Arduino\\Uno R3};
    \node[font=\sffamily\scriptsize, text=gray!70!black] at (2.0, -0.6) {(ATmega328P)};

    % Alimentacao do Arduino
    \draw (0, 5.8) -- (-0.8, 5.8) node[vcc]{+5V};
    \node[right, font=\ttfamily\footnotesize] at (0, 5.8) {5V};

    \draw (0, -5.8) -- (-0.8, -5.8) node[ground]{};
    \node[right, font=\ttfamily\footnotesize] at (0, -5.8) {GND};

    % =========================================================================
    % 2. BLOCO DIREITO: 8 LEDs (LED0 a LED7) + LED_BUILTIN + RESISTORES DE 220 OHMS
    % =========================================================================
    \foreach \y/\pin/\dname/\rname in {
        6.0/D2/LED0/R0,
        4.5/D3/LED1/R1,
        3.0/D4/LED2/R2,
        1.5/D5/LED3/R3,
        0.0/D6/LED4/R4,
        -1.5/D7/LED5/R5,
        -3.0/D8/LED6/R6,
        -4.5/D9/LED7/R7,
        -6.0/D13/{LED\_BUILTIN}/R8%
    }{
        \node[left, font=\ttfamily\footnotesize] at (4.0, \y) {\pin};
        \draw (4.0, \y) -- (4.8, \y)
            to[leDo, l^={\dname}] (7.2, \y)
            to[R, l^=\rname, a_=220<\ohm>] (9.8, \y)
            -- (10.6, \y);
    }

    % Barramento comum de GND (sem ponto no canto superior de 90 graus)
    \draw (10.6, 6.0) -- (10.6, -7.0) node[ground]{};
    \foreach \y in {4.5, 3.0, 1.5, 0.0, -1.5, -3.0, -4.5, -6.0}{
        \node[circ] at (10.6, \y) {};
    }

\end{circuitikz}
\end{document}